\documentclass[11pt,a4paper]{article}
\usepackage[utf8]{inputenc}
\usepackage[T1]{fontenc}
\usepackage{amsmath,amssymb,amsfonts}
\usepackage{physics}
\usepackage{booktabs}
\usepackage{enumitem}
\usepackage[colorlinks=true,linkcolor=blue,citecolor=blue,urlcolor=blue]{hyperref}
\usepackage[margin=1in]{geometry}
\usepackage{parskip}

\title{Numerical Wave Packet Propagation for the Q-Ontic Lab\\Double-Slit Simulator: Analysis \& Proposal}
\author{Q-Ontic Lab Development}
\date{February 2026}

\begin{document}
\maketitle

\section{Current Implementation}

The Q-Ontic Lab double-slit simulator currently uses a \textbf{fully analytical Huygens--Fresnel point-source model}. Each slit emits a 2D cylindrical wavelet:
\begin{equation}
  \psi_{\text{point}}(\mathbf{r}, t)
  = \frac{A}{r}\,e^{i(kr - \omega t)},
  \qquad r = \sqrt{x^2 + y^2},
\end{equation}
and the total wave function after the barrier is the coherent superposition of the wavelets from each open slit.

\subsection{Key characteristics}
\begin{itemize}
  \item \textbf{Monochromatic, steady-state, continuous-wave} --- not a wave packet.
  \item Time evolution is a trivial global phase rotation $e^{-i\omega t}$, allowing the base $\psi$ at $t=0$ to be cached and rotated.
  \item The ``Wave Packet Width'' UI parameter is a \emph{display mask} that restricts visibility to a radial band around particle positions, creating a visual packet effect without changing the underlying physics.
  \item Bohmian velocity is computed analytically:
    \begin{equation}
      \mathbf{v} = \frac{\hbar}{m}\,
      \frac{\operatorname{Im}(\psi^*\,\nabla\psi)}{|\psi|^2}.
    \end{equation}
  \item The quantum potential is computed via finite-difference Laplacian of $R = |\psi|$.
  \item Dispersion relation: photon $\omega = kc$; massive particle $\omega = \hbar k^2 / 2m$.
\end{itemize}

\subsection{Existing infrastructure}
\begin{table}[h]
\centering
\begin{tabular}{@{}ll@{}}
\toprule
\textbf{Resource} & \textbf{Details} \\
\midrule
WebGL 1.0 & Full fragment shader pipeline; per-pixel wave rendering at 60\,fps \\
Web Worker & Offloads CDF precomputation; extensible for PDE integration \\
Float32Arrays & \texttt{basePsiReal}, \texttt{basePsiImag}, \texttt{basePsiPhase} per-pixel buffers \\
Frame cache & 64 frames/period with LRU eviction \\
CPU fallback & \texttt{generateWaveInfo()} loops over every pixel --- directly replaceable \\
\bottomrule
\end{tabular}
\end{table}


\section{Numerical Methods for Wave Packet Propagation}

Below are the principal approaches for replacing the analytical Huygens wavelet with a time-dependent Schr\"odinger equation (TDSE) solver that propagates a localized Gaussian wave packet through a double-slit potential.

\subsection{Split-Step Fourier Method (SSFM)}

Alternate between position-space potential kicks and momentum-space kinetic kicks each $\Delta t$:
\begin{equation}
  \psi(t + \Delta t)
  = e^{-iV\Delta t/2\hbar}\;
    \mathcal{F}^{-1}\!\Bigl[
      e^{-i\hbar|\mathbf{k}|^2\Delta t/2m}\;
      \mathcal{F}\!\bigl[
        e^{-iV\Delta t/2\hbar}\,\psi(t)
      \bigr]
    \Bigr].
\end{equation}

\begin{itemize}
  \item \textbf{Pros:} Spectrally accurate for the kinetic operator; $O(N\log N)$ per step via FFT; straightforward implementation.
  \item \textbf{Cons:} Requires FFT (JS library or WebGL/WebGPU compute); periodic boundary conditions necessitate absorbing layers; sharp potentials can cause Gibbs ringing.
  \item \textbf{Browser fit:} Good. JS FFT libraries exist (e.g., \texttt{fft.js}). Also implementable on GPU via WebGL ping-pong textures with FFT shaders, or via WebGPU compute.
\end{itemize}

\subsection{Crank--Nicolson Finite Difference}

Implicit trapezoidal time integration of the discretized Hamiltonian:
\begin{equation}
  \Bigl(1 + \frac{i\Delta t}{2\hbar}\hat{H}\Bigr)\psi^{n+1}
  = \Bigl(1 - \frac{i\Delta t}{2\hbar}\hat{H}\Bigr)\psi^n.
\end{equation}

\begin{itemize}
  \item \textbf{Pros:} Unconditionally stable; exactly unitary (norm-conserving); second-order accurate in space and time.
  \item \textbf{Cons:} In 2D requires ADI (Alternating Direction Implicit) splitting into tridiagonal systems; larger implementation complexity.
  \item \textbf{Browser fit:} Good for CPU via Web Workers. Each 1D tridiagonal system is $O(N)$ via the Thomas algorithm. Row/column sweeps are embarrassingly parallel.
\end{itemize}

\subsection{Explicit Leapfrog FDTD}

Evolve real and imaginary parts on staggered time grids using a 5-point Laplacian stencil:
\begin{align}
  \operatorname{Re}(\psi^{n+1})
    &= \operatorname{Re}(\psi^{n-1})
     + \frac{2\Delta t}{\hbar}\,\hat{H}\,\operatorname{Im}(\psi^n), \\[4pt]
  \operatorname{Im}(\psi^{n+1})
    &= \operatorname{Im}(\psi^{n-1})
     - \frac{2\Delta t}{\hbar}\,\hat{H}\,\operatorname{Re}(\psi^n).
\end{align}

\begin{itemize}
  \item \textbf{Pros:} Very simple; no linear algebra; trivially parallelizable stencil operation; maps directly to GPU fragment shaders via ping-pong textures.
  \item \textbf{Cons:} Conditionally stable: $\Delta t \le m\,\Delta x^2 / (\hbar\,d)$ where $d$ is dimensionality. Not exactly unitary (slow norm drift without symplectic correction).
  \item \textbf{Browser fit:} Excellent for WebGL. A two-texture ping-pong with a 5-point stencil fragment shader can run at 60\,fps for $512\times512$ grids.
\end{itemize}

\subsection{Kirchhoff--Fresnel Integral Upgrade (Semi-Analytical)}

Upgrade the current point-source model to integrate the Huygens--Fresnel--Kirchhoff diffraction integral across each slit's finite aperture:
\begin{equation}
  \psi(P) = -\frac{i}{\lambda}
  \int_{\text{slit}} \psi_{\text{inc}}(Q)\,
  \frac{e^{ikr}}{r}\,K(\chi)\,dS,
\end{equation}
where $K(\chi)$ is the obliquity factor. For a wave packet, convolve with the packet's spectral profile (integration over $\Delta k$).

\begin{itemize}
  \item \textbf{Pros:} Keeps existing architecture mostly intact; handles finite slit width; can incorporate wave packet envelopes via spectral convolution.
  \item \textbf{Cons:} Not a true PDE solver; less flexible for complex potentials; limited Bohmian trajectory support beyond the far field.
  \item \textbf{Browser fit:} Excellent. The slit integral is 1D and cheap.
\end{itemize}

\subsection{Hybrid: Analytical Propagation + Numerical Slit Interaction}

Use exact free-particle Gaussian wave packet propagation in the regions before and after the barrier:
\begin{equation}
  \psi(x,t) = \biggl(\frac{2\alpha}{\pi}\biggr)^{\!1/4}
  \frac{1}{\sqrt{1 + 2i\alpha\hbar t/m}}\,
  \exp\!\biggl(
    -\frac{\alpha(x - x_0 - v_0 t)^2}{1 + 2i\alpha\hbar t/m}
    + ik_0 x - i\omega_0 t
  \biggr),
\end{equation}
and solve numerically only in a narrow strip around the slit wall.

\begin{itemize}
  \item \textbf{Pros:} Minimal computational cost; exact in free regions; analytical dispersion.
  \item \textbf{Cons:} Boundary matching at the numerical/analytical interface requires care; limited to simple potentials.
\end{itemize}


\section{Recommendation}

Given the existing WebGL pipeline and real-time performance requirements, the recommended implementation path is:

\begin{enumerate}
  \item \textbf{Phase 1 --- Explicit Leapfrog FDTD on WebGL} (primary recommendation)
    \begin{itemize}
      \item Add a two-texture ping-pong pass before the display shader, storing $\operatorname{Re}(\psi)$ and $\operatorname{Im}(\psi)$ in float textures.
      \item Initial condition: 2D Gaussian wave packet
        $\psi_0(x,y) = \exp\bigl[-\alpha\bigl((x-x_0)^2 + (y-y_0)^2\bigr)\bigr]\,e^{ik_0 x}$.
      \item Barrier potential: $V(x,y) = V_0$ inside the wall material, $V = 0$ elsewhere.
      \item Absorbing boundaries: imaginary potential ramp at domain edges to prevent reflections.
      \item All existing Bohmian mechanics code works unchanged --- $\nabla\psi$ computed from the grid via finite differences (already done for the quantum potential).
    \end{itemize}

  \item \textbf{Phase 2 --- Split-Step FFT upgrade} (optional, for improved accuracy)
    \begin{itemize}
      \item Run in a Web Worker using a JS FFT library, writing results into existing \texttt{basePsiReal}/\texttt{basePsiImag} buffers, then upload to WebGL textures.
      \item Gives spectral accuracy and correct dispersion handling.
    \end{itemize}
\end{enumerate}

\subsection{Design principle: switchable modes}

The implementation should allow toggling between the current analytical Huygens--Fresnel mode and the new numerical wave packet mode at runtime, via a UI control. Both modes share the same rendering, Bohmian trajectory, and particle detection pipelines --- only the wave function source differs.

\subsection{Physics gained}

A true numerical wave packet simulation would add:
\begin{itemize}
  \item \textbf{Dispersion:} the packet spreads over time as $\sigma(t) = \sigma_0\sqrt{1 + (\hbar t / 2m\sigma_0^2)^2}$.
  \item \textbf{Finite slit width diffraction:} proper single-slit diffraction envelope modulating the double-slit interference.
  \item \textbf{Temporal dynamics:} the wave arrives, passes through the slits, and interferes --- a transient process rather than a steady state.
  \item \textbf{Realistic barrier interaction:} partial reflection, evanescent tunneling through narrow barriers.
\end{itemize}

\subsection{Key references}

\begin{itemize}
  \item A.~Goldberg, H.\,M.~Schey, and J.\,L.~Schwartz, ``Computer-generated motion pictures of one-dimensional quantum-mechanical transmission and reflection phenomena,'' \emph{Am.\ J.\ Phys.}\ \textbf{35}, 177 (1967).
  \item M.\,D.~Feit, J.\,A.~Fleck, and A.~Steiger, ``Solution of the Schr\"odinger equation by a spectral method,'' \emph{J.\ Comput.\ Phys.}\ \textbf{47}, 412 (1982). [Split-step method]
  \item R.~Kosloff, ``Time-dependent quantum-mechanical methods for molecular dynamics,'' \emph{J.\ Phys.\ Chem.}\ \textbf{92}, 2087 (1988).
  \item A.~Askar and A.\,S.~Cakmak, ``Explicit integration method for the time-dependent Schr\"odinger equation for collision problems,'' \emph{J.\ Chem.\ Phys.}\ \textbf{68}, 2794 (1978). [Leapfrog scheme]
\end{itemize}

\end{document}
